\section{Mathematische Analyse -- Formelsammlung}

Diese Section sammelt die \textbf{wichtigsten mathematischen Werkzeuge}
für die Analyse in der Leistungselektronik:
\textbf{Ableitungen}, \textbf{Integrale}, \textbf{Trigonometrie},
\textbf{Exponentialfunktionen}, \textbf{Komplexrechnung},
\textbf{Laplace-Transformation} und \textbf{Differentialgleichungen (RLC)}.
Sie ist als \textbf{reine Nachschlage-Section} gedacht.

\smallskip

\textbf{Ziel:}
\begin{itemize}
    \item Schneller Zugriff auf Standardformeln
    \item Fehlerfreie Umformungen in Prüfungen
    \item Einheitliche Notation im ganzen Skript
\end{itemize}

% -------------------------------------------------
\subsection{Grundlegende Ableitungsregeln}

\subsubsection{Elementare Funktionen}

\[
\frac{d}{dx} c = 0
\qquad
\frac{d}{dx} x = 1
\qquad
\frac{d}{dx} x^n = n x^{n-1}
\]

\[
\frac{d}{dx} \e^{x} = \e^{x}
\qquad
\frac{d}{dx} \ln x = \frac{1}{x}
\]

\subsubsection{Lineare Regeln}

\[
\frac{d}{dx} [f(x) + g(x)] = f'(x) + g'(x)
\qquad
\frac{d}{dx} [c \, f(x)] = c \, f'(x)
\]

\subsubsection{Produktregel}

\[
\cbl{\frac{d}{dx} [f(x) g(x)] = f'(x) g(x) + f(x) g'(x)}
\]

\subsubsection{Quotientenregel}

\[
\cbl{\frac{d}{dx} \left[ \frac{f(x)}{g(x)} \right]
= \frac{f'(x) g(x) - f(x) g'(x)}{g(x)^2}}
\]

\subsubsection{Kettenregel}

\[
\cbl{\frac{d}{dx} f(g(x)) = f'(g(x)) \cdot g'(x)}
\]

% -------------------------------------------------
\subsection{Ableitungen trigonometrischer Funktionen}

\[
\cbl{\frac{d}{dx} \sin x = \cos x}
\qquad
\cbl{\frac{d}{dx} \cos x = - \sin x}
\]

\[
\frac{d}{dx} \tan x = \frac{1}{\cos^2 x}
\qquad
\frac{d}{dx} \cot x = - \frac{1}{\sin^2 x}
\]

\[
\frac{d}{dx} \arcsin x = \frac{1}{\sqrt{1-x^2}}
\qquad
\frac{d}{dx} \arccos x = -\frac{1}{\sqrt{1-x^2}}
\]

% -------------------------------------------------
\subsection{Grundlegende Integralregeln}

\subsubsection{Elementare Integrale}

\[
\int c \, dx = c x + C
\qquad
\int x^n \, dx = \frac{x^{n+1}}{n+1} + C \quad (n \neq -1)
\]

\[
\int \frac{1}{x} \, dx = \ln |x| + C
\qquad
\int \e^{x} \, dx = \e^{x} + C
\]

\subsubsection{Lineare Regeln}

\[
\int [f(x) + g(x)] dx = \int f(x) dx + \int g(x) dx
\qquad
\int c f(x) dx = c \int f(x) dx
\]

% -------------------------------------------------
\subsection{Integrale trigonometrischer Funktionen}

\[
\cbl{\int \sin x \, dx = - \cos x + C}
\qquad
\cbl{\int \cos x \, dx = \sin x + C}
\]

\[
\int \tan x \, dx = - \ln |\cos x| + C
\qquad
\int \cot x \, dx = \ln |\sin x| + C
\]

\[
\int \sin^2 x \, dx = \frac{x}{2} - \frac{\sin(2x)}{4} + C
\qquad
\int \cos^2 x \, dx = \frac{x}{2} + \frac{\sin(2x)}{4} + C
\]

\[
\int \sin(ax) \, dx = -\frac{1}{a}\cos(ax) + C
\qquad
\int \cos(ax) \, dx = \frac{1}{a}\sin(ax) + C
\]

% -------------------------------------------------
\subsection{Partielle Integration}

\[
\cbl{\int u \, dv = u v - \int v \, du}
\]

Beispiel:
\[
\int x \e^{x} dx = x \e^{x} - \int \e^{x} dx = (x-1) \e^{x} + C
\]

% -------------------------------------------------
\subsection{Substitution}

\[
\cbl{u = g(x), \qquad du = g'(x) dx}
\]

Beispiel:
\[
\int \sin(3x) dx
= -\frac{1}{3} \cos(3x) + C
\]

% -------------------------------------------------
\subsection{Trigonometrische Identitäten}

\subsubsection{Grundidentitäten}

\[
\cbl{\sin^2 x + \cos^2 x = 1}
\qquad
\cbl{1 + \tan^2 x = \frac{1}{\cos^2 x}}
\qquad
1 + \cot^2 x = \frac{1}{\sin^2 x}
\]

\subsubsection{Winkeladdition}

\[
\cbl{\sin(a \pm b) = \sin a \cos b \pm \cos a \sin b}
\]

\[
\cbl{\cos(a \pm b) = \cos a \cos b \mp \sin a \sin b}
\]

\subsubsection{Doppelwinkel}

\[
\cbl{\sin(2x) = 2 \sin x \cos x}
\qquad
\cbl{\cos(2x) = \cos^2 x - \sin^2 x = 1 - 2 \sin^2 x = 2 \cos^2 x - 1}
\]

\subsubsection{Halbwinkel}

\[
\cbl{\sin^2\left(\frac{x}{2}\right) = \frac{1-\cos x}{2}}
\qquad
\cbl{\cos^2\left(\frac{x}{2}\right) = \frac{1+\cos x}{2}}
\]

% -------------------------------------------------
\subsection{Produkt--Summen--Formeln}

\[
\cbl{\sin a \cos b = \frac{1}{2} [ \sin(a+b) + \sin(a-b) ]}
\]

\[
\cbl{\cos a \cos b = \frac{1}{2} [ \cos(a+b) + \cos(a-b) ]}
\]

\[
\cbl{\sin a \sin b = \frac{1}{2} [ \cos(a-b) - \cos(a+b) ]}
\]

\[
\cbl{\sin a + \sin b = 2 \sin\left(\frac{a+b}{2}\right) \cos\left(\frac{a-b}{2}\right)}
\]

\[
\cbl{\cos a + \cos b = 2 \cos\left(\frac{a+b}{2}\right) \cos\left(\frac{a-b}{2}\right)}
\]

\[
\cbl{\cos a - \cos b = -2 \sin\left(\frac{a+b}{2}\right) \sin\left(\frac{a-b}{2}\right)}
\]

% -------------------------------------------------
\subsection{Exponential- und Euler-Formeln}

\subsubsection{Euler-Formel}

\[
\cbl{\e^{\jimg x} = \cos x + \jimg \sin x}
\qquad
\cbl{\e^{-\jimg x} = \cos x - \jimg \sin x}
\]

\[
\cbl{\cos x = \frac{\e^{\jimg x} + \e^{-\jimg x}}{2}}
\qquad
\cbl{\sin x = \frac{\e^{\jimg x} - \e^{-\jimg x}}{2\jimg}}
\]

\subsubsection{Komplexe Sinusdarstellung}

\[
\cbl{\hat X \sin(\omega t + \varphi) = \Im\left\{ \hat X \e^{\jimg(\omega t + \varphi)} \right\}}
\]

\[
\cbl{\hat X \cos(\omega t + \varphi) = \Re\left\{ \hat X \e^{\jimg(\omega t + \varphi)} \right\}}
\]

% -------------------------------------------------
\subsection{Komplexrechnung}

\subsubsection{Grundformen}

\[
\cbl{z = x + \jimg y}
\qquad
\cbl{z = r \e^{\jimg \varphi}}
\]

\[
r = |z| = \sqrt{x^2 + y^2}
\qquad
\varphi = \arg(z)
\]

\subsubsection{Konjugation}

\[
\cbl{\overline{z} = x - \jimg y}
\qquad
\cbl{z \overline{z} = |z|^2}
\]

\subsubsection{Multiplikation und Division}

\[
\cbl{z_1 z_2 = r_1 r_2 \e^{\jimg (\varphi_1 + \varphi_2)}}
\qquad
\cbl{\frac{z_1}{z_2} = \frac{r_1}{r_2} \e^{\jimg (\varphi_1 - \varphi_2)}}
\]

\subsubsection{Real- und Imaginärteil}

\[
\cbl{\Re(z) = x}
\qquad
\cbl{\Im(z) = y}
\]

% -------------------------------------------------
\subsection{Komplexe Wechselstromrechnung}

\subsubsection{Ableiten im Zeitbereich}

\[
\cbl{\frac{d}{dt} \e^{\jimg \omega t} = \jimg \omega \e^{\jimg \omega t}}
\]

\subsubsection{Impedanzen}

\[
\cbl{Z_R = R}
\qquad
\cbl{Z_L = \jimg \omega L}
\qquad
\cbl{Z_C = \frac{1}{\jimg \omega C}}
\]

\subsubsection{Serien- und Parallelschaltung}

\[
\cbl{Z_{\mathrm{Serie}} = Z_1 + Z_2 + \dots}
\]

\[
\cbl{\frac{1}{Z_{\mathrm{Par}}} = \frac{1}{Z_1} + \frac{1}{Z_2} + \dots}
\]

% -------------------------------------------------
\subsection{Orthogonalitätsintegrale (Fourier-Kern)}

\[
\cbl{\int_0^T \cos(m \omega t)\cos(n \omega t)\,dt =
\begin{cases}
T, & m=n=0 \\
\frac{T}{2}, & m=n>0 \\
0, & m\neq n
\end{cases}}
\]

\[
\cbl{\int_0^T \sin(m \omega t)\sin(n \omega t)\,dt =
\begin{cases}
\frac{T}{2}, & m=n \\
0, & m\neq n
\end{cases}}
\]

\[
\cbl{\int_0^T \sin(m \omega t)\cos(n \omega t)\,dt = 0}
\]

% -------------------------------------------------
\subsection{Grenzwerte und Reihen}

\[
\cbl{\lim_{x \to 0} \frac{\sin x}{x} = 1}
\qquad
\cbl{\lim_{x \to 0} \frac{1 - \cos x}{x^2} = \frac{1}{2}}
\]

\[
\cbl{\e^{x} = \sum_{n=0}^{\infty} \frac{x^n}{n!}}
\qquad
\cbl{\sin x = \sum_{n=0}^{\infty} (-1)^n \frac{x^{2n+1}}{(2n+1)!}}
\qquad
\cbl{\cos x = \sum_{n=0}^{\infty} (-1)^n \frac{x^{2n}}{(2n)!}}
\]

% -------------------------------------------------
\subsection{Bestimmte Integrale über Perioden}

\[
\cbl{\int_0^T \sin(n \omega t)\,dt = 0}
\qquad
\cbl{\int_0^T \cos(n \omega t)\,dt = 0}
\]

\[
\cbl{\int_0^T \sin^2(n \omega t)\,dt = \frac{T}{2}}
\qquad
\cbl{\int_0^T \cos^2(n \omega t)\,dt = \frac{T}{2}}
\]

% =================================================
\subsection{Laplace-Transformation}

\subsubsection{Definition}

Für eine zeitkontinuierliche Funktion $f(t)$ (typisch $t\ge 0$):
\[
\cbl{\mathcal{L}\{f(t)\} = F(s) = \int_{0}^{\infty} f(t)\, \e^{-st}\,dt}
\qquad
s = \sigma + \jimg \omega
\]

\subsubsection{Wichtige Grundtransformierte}

\[
\mathcal{L}\{1\} = \frac{1}{s}
\qquad
\mathcal{L}\{t\} = \frac{1}{s^2}
\qquad
\mathcal{L}\{t^n\} = \frac{n!}{s^{n+1}}
\]

\[
\cbl{\mathcal{L}\{\e^{at}\} = \frac{1}{s-a}}
\]

\[
\cbl{\mathcal{L}\{\sin(\omega t)\} = \frac{\omega}{s^2+\omega^2}}
\qquad
\cbl{\mathcal{L}\{\cos(\omega t)\} = \frac{s}{s^2+\omega^2}}
\]

\[
\mathcal{L}\{\sinh(\omega t)\} = \frac{\omega}{s^2-\omega^2}
\qquad
\mathcal{L}\{\cosh(\omega t)\} = \frac{s}{s^2-\omega^2}
\]

\subsubsection{Linearität}

\[
\cbl{\mathcal{L}\{a f(t) + b g(t)\} = aF(s) + bG(s)}
\]

\subsubsection{Ableitungssatz}

\[
\cbl{\mathcal{L}\left\{\frac{df}{dt}\right\} = sF(s) - f(0)}
\]

\[
\cbl{\mathcal{L}\left\{\frac{d^2f}{dt^2}\right\} = s^2F(s) - s f(0) - f'(0)}
\]

\crd{\textrightarrow\ Anfangswerte gehen als Terme in den Zähler. Häufige Prüfungsfalle.}

\subsubsection{Integralsatz}

\[
\cbl{\mathcal{L}\left\{\int_0^t f(\tau)\,d\tau\right\} = \frac{1}{s} F(s)}
\]

\subsubsection{Zeitverschiebung (Heaviside)}

\[
\cbl{\mathcal{L}\{f(t-t_0)u(t-t_0)\} = \e^{-s t_0} F(s)}
\]

\subsubsection{Frequenzverschiebung}

\[
\cbl{\mathcal{L}\{\e^{at} f(t)\} = F(s-a)}
\]

\subsubsection{Faltung}

\[
\cbl{(f*g)(t) = \int_0^t f(\tau)g(t-\tau)\,d\tau}
\qquad
\cbl{\mathcal{L}\{f*g\} = F(s)G(s)}
\]

\subsubsection{Initial- und Finalwertsatz}

\[
\cbl{f(0^+) = \lim_{s\to\infty} sF(s)}
\qquad
\cbl{f(\infty) = \lim_{s\to 0} sF(s)}
\]

\crd{\textrightarrow\ Finalwertsatz nur gültig, wenn alle Pole von $sF(s)$ in der linken Halbebene liegen.}

\subsubsection{Inverse Laplace (Standardstrategien)}

\begin{itemize}
    \item Partialbruchzerlegung
    \item Tabellenrücktransformation
    \item Verschiebungssätze (Zeit, Frequenz)
    \item Quadratische Ergänzung für $s^2+\omega^2$
\end{itemize}

% =================================================
\subsection{Übertragungsfunktionen und Systemdarstellung}

Für lineare zeitinvariante Systeme (LTI):
\[
\cbl{H(s) = \frac{Y(s)}{X(s)}}
\]

\subsubsection{Pol-Nullstellen-Darstellung}

\[
\cbl{H(s)=K \frac{(s-z_1)(s-z_2)\dots}{(s-p_1)(s-p_2)\dots}}
\]

\subsubsection{Zeitkonstante und Eckfrequenz}

Erste Ordnung:
\[
\cbl{H(s)=\frac{1}{1+s\tau}}
\qquad
\cbl{\omega_c = \frac{1}{\tau}}
\qquad
\cbl{f_c = \frac{1}{2\pi\tau}}
\]

% =================================================
\subsection{Differentialgleichungen 1. Ordnung (R--L / R--C)}

\subsubsection{Standardform}

\[
\cbl{\tau \frac{dy}{dt} + y = K \, u(t)}
\]

Homogene Lösung:
\[
\cbl{y_h(t) = C \e^{-t/\tau}}
\]

Partikuläre Lösung (Konstantanregung $u(t)=U$):
\[
\cbl{y_p(t)=K U}
\]

Gesamtlösung:
\[
\cbl{y(t) = K U + \big(y(0)-K U\big)\e^{-t/\tau}}
\]

\subsubsection{R--L-Kreis}

\[
\cbl{L\frac{di}{dt} + Ri = u(t)}
\qquad
\cbl{\tau = \frac{L}{R}}
\]

Sprung $u(t)=U$:
\[
\cbl{i(t) = \frac{U}{R}\left(1-\e^{-t/\tau}\right) + i(0)\e^{-t/\tau}}
\]

\subsubsection{R--C-Kreis}

\[
\cbl{RC\frac{du_C}{dt} + u_C = u(t)}
\qquad
\cbl{\tau = RC}
\]

Sprung $u(t)=U$:
\[
\cbl{u_C(t) = U\left(1-\e^{-t/\tau}\right) + u_C(0)\e^{-t/\tau}}
\]

\crd{\textrightarrow\ Bei Induktivitäten ist $i(t)$ stetig, bei Kapazitäten ist $u_C(t)$ stetig.}

% =================================================
\subsection{Differentialgleichungen 2. Ordnung (RLC-Systeme)}

\subsubsection{Standardform}

\[
\cbl{\frac{d^2y}{dt^2} + 2\zeta\omega_0 \frac{dy}{dt} + \omega_0^2 y = K \omega_0^2 u(t)}
\]

\[
\cbl{\omega_0 = \sqrt{\frac{1}{LC}}}
\qquad
\cbl{\zeta = \frac{R}{2}\sqrt{\frac{C}{L}}}
\]

\subsubsection{Dämpfungsfälle}

\begin{itemize}
    \item \textbf{Unterdämpft:} $0<\zeta<1$ (schwingend)
    \item \textbf{Kritisch gedämpft:} $\zeta=1$
    \item \textbf{Überdämpft:} $\zeta>1$
\end{itemize}

Unterdämpfte Eigenfrequenz:
\[
\cbl{\omega_d = \omega_0\sqrt{1-\zeta^2}}
\]

\subsubsection{Freie Schwingung (homogene Lösung)}

Unterdämpft:
\[
\cbl{y_h(t)=\e^{-\zeta\omega_0 t}\left(C_1\cos(\omega_d t)+C_2\sin(\omega_d t)\right)}
\]

Überdämpft:
\[
\cbl{y_h(t)=C_1\e^{s_1 t}+C_2\e^{s_2 t}}
\qquad
s_{1,2}=-\omega_0\left(\zeta\pm\sqrt{\zeta^2-1}\right)
\]

Kritisch:
\[
\cbl{y_h(t)=(C_1 + C_2 t)\e^{-\omega_0 t}}
\]

\subsubsection{Serien-RLC (Stromdynamik)}

KVL:
\[
\cbl{u(t)=Ri(t)+L\frac{di}{dt}+\frac{1}{C}\int i(t)\,dt}
\]

Differenziert:
\[
\cbl{\frac{d^2 i}{dt^2} + \frac{R}{L}\frac{di}{dt} + \frac{1}{LC}i = \frac{1}{L}\frac{du}{dt}}
\]

\subsubsection{Parallele-RLC (Spannungsdynamik)}

KCL:
\[
i(t)=\frac{u(t)}{R}+C\frac{du}{dt}+\frac{1}{L}\int u(t)\,dt
\]

Differenziert:
\[
\cbl{\frac{d^2 u}{dt^2} + \frac{1}{RC}\frac{du}{dt} + \frac{1}{LC}u = \frac{1}{C}\frac{di}{dt}}
\]

\subsubsection{Laplace-Lösung von DGLs}

Vorgehen:
\begin{enumerate}
    \item DGL aufstellen (KVL/KCL)
    \item Laplace anwenden (Ableitungssätze)
    \item Nach $Y(s)$ auflösen
    \item Partialbrüche
    \item Rücktransformation
\end{enumerate}

\subsubsection{Standardübertragungsfunktionen}

\paragraph{Tiefpass 1. Ordnung}
\[
\cbl{H(s)=\frac{1}{1+s\tau}}
\]

\paragraph{Bandpass / Resonanz 2. Ordnung (Normalform)}
\[
\cbl{H(s)=\frac{\omega_0^2}{s^2+2\zeta\omega_0 s+\omega_0^2}}
\]

% -------------------------------------------------
\subsection{Typische Merksätze (Analyse-Kern)}

\begin{itemize}
    \item Kettenregel, Produktregel, trigonometrische Identitäten sicher beherrschen.
    \item Laplace macht aus DGLs algebraische Gleichungen.
    \item $i_L(t)$ und $u_C(t)$ sind stetig.
    \item Polstellen bestimmen Stabilität und Zeitverhalten.
\end{itemize}

\subsection{Typische Fehlerquellen}

\begin{itemize}
    \item Anfangswerte bei Laplace-Ableitung vergessen.
    \item Finalwertsatz ohne Polprüfung verwendet.
    \item Vorzeichenfehler bei $\cos'(x)$ und bei KVL/KCL.
    \item $u_C$ bzw. $i_L$ als sprungfähig angenommen.
\end{itemize}
